When multilingual large language models encounter conflicting information between languages, presentation language can determine the answer, but the internal mechanism remains unknown. We probe residual streams in Llama-3.1-8B-Instruct, Qwen3-8B, and Gemma-2-9B-IT on English--Chinese conflicts, combining cross-lingual transfer, Cross-lingual Natural Language Inference (XNLI) non-conflict baselines, confound exclusion, and representation steering. Conflict signals are linearly detectable across both languages, and cross-lingual transfer exhibits a directional asymmetry: Chinese-trained probes transfer to English ${\sim}$7 percentage points more accurately than the reverse, matching the XNLI baseline and consistent with general multilingual structure rather than conflict-specific encoding. A dissociation emerges between detection and decision: the probe signal is broadly distributed across mid-to-upper layers, yet representation steering achieves causal influence only at a mid-layer (L13, 41\% depth) in Llama; analogous Qwen and Gemma steering yields no significant effect. These findings indicate that linear detectability does not imply causal relevance, with implications for probe-based intervention\footnote{Code and data are available at \url{https://anonymous.4open.science/r/clkc_mech_interp-F894/}.}.
